\documentclass[10pt,a4paper]{amsart}
\usepackage{amsmath,amssymb,amsthm,graphicx}
% Retained preamble for the Stacks Project chapter snapshot in this directory.
% The chapter's own preamble.tex is a copy of the shared Stacks Project preamble; the block
% below is sliced MECHANICALLY, line for line, from it. Nothing was added, removed or
% redefined; the wrapper class, the xr-hyper/hyperref wiring and the externaldocument list
% are omitted because the assembler supplies the document class.
%
% lines taken from preamble.tex: 3, 5, 9, 11, 12, 156, 157, 158, ... (51 lines total)
%
% package coverage restored from the chapter's own preamble (2026-10-13): xy, xy-2cell

\usepackage[all]{xy}
\xyoption{2cell}
\UseAllTwocells

\newenvironment{reference}{\comment}{\endcomment}
\newenvironment{slogan}{\comment}{\endcomment}
\newenvironment{history}{\comment}{\endcomment}
\theoremstyle{plain}
\newtheorem{theorem}[subsection]{Theorem}
\newtheorem{proposition}[subsection]{Proposition}
\newtheorem{lemma}[subsection]{Lemma}
\theoremstyle{definition}
\newtheorem{definition}[subsection]{Definition}
\newtheorem{example}[subsection]{Example}
\newtheorem{exercise}[subsection]{Exercise}
\newtheorem{situation}[subsection]{Situation}
\theoremstyle{remark}
\newtheorem{remark}[subsection]{Remark}
\newtheorem{remarks}[subsection]{Remarks}
\numberwithin{equation}{subsection}
\def\lim{\mathop{\mathrm{lim}}\nolimits}
\def\colim{\mathop{\mathrm{colim}}\nolimits}
\def\Spec{\mathop{\mathrm{Spec}}}
\def\Hom{\mathop{\mathrm{Hom}}\nolimits}
\def\Ext{\mathop{\mathrm{Ext}}\nolimits}
\def\SheafHom{\mathop{\mathcal{H}\!\mathit{om}}\nolimits}
\def\SheafExt{\mathop{\mathcal{E}\!\mathit{xt}}\nolimits}
\def\Sch{\mathit{Sch}}
\def\Mor{\mathop{\mathrm{Mor}}\nolimits}
\def\Ob{\mathop{\mathrm{Ob}}\nolimits}
\def\Sh{\mathop{\mathit{Sh}}\nolimits}
\def\NL{\mathop{N\!L}\nolimits}
\def\CH{\mathop{\mathrm{CH}}\nolimits}
\def\proetale{{pro\text{-}\acute{e}tale}}
\def\etale{{\acute{e}tale}}
\def\QCoh{\mathit{QCoh}}
\def\Ker{\mathop{\mathrm{Ker}}}
\def\Im{\mathop{\mathrm{Im}}}
\def\Coker{\mathop{\mathrm{Coker}}}
\def\Coim{\mathop{\mathrm{Coim}}}
\DeclareMathSymbol{\boxtimes}{\mathbin}{AMSa}{"02}
\def\QCohstack{\mathcal{QC}\!\mathit{oh}}
\def\Cohstack{\mathcal{C}\!\mathit{oh}}
\def\Spacesstack{\mathcal{S}\!\mathit{paces}}
\def\Quotfunctor{\mathrm{Quot}}
\def\Hilbfunctor{\mathrm{Hilb}}
\def\Curvesstack{\mathcal{C}\!\mathit{urves}}
\def\Polarizedstack{\mathcal{P}\!\mathit{olarized}}
\def\Complexesstack{\mathcal{C}\!\mathit{omplexes}}
\def\Pic{\mathop{\mathrm{Pic}}\nolimits}
\def\Picardstack{\mathcal{P}\!\mathit{ic}}
\def\Picardfunctor{\mathrm{Pic}}
\def\Deformationcategory{\mathcal{D}\!\mathit{ef}}
\usepackage{cleveref}
\crefname{originallemma}{original lemma}{original lemmas}
\Crefname{originallemma}{Original lemma}{Original lemmas}
\makeatletter
\newlabel{definition-hypercovering}{{01G5}{source}{}{sourcepointer-1}{}}
\newlabel{definition-hypercovering@cref}{{[originallemma][0][]01G5}{[1][0][]source}{}{}{}}
\makeatother
\crefname{originallemma}{original lemma}{original lemmas}
\Crefname{originallemma}{Original lemma}{Original lemmas}
\makeatletter
\newlabel{lemma-covering-permanence}{{01G4}{source}{}{sourcepointer-2}{}}
\newlabel{lemma-covering-permanence@cref}{{[originallemma][0][]01G4}{[1][0][]source}{}{}{}}
\makeatother
\crefname{originallemma}{original lemma}{original lemmas}
\Crefname{originallemma}{Original lemma}{Original lemmas}
\makeatletter
\newlabel{simplicial-lemma-cosk-fibre-product}{{018D}{source}{}{sourcepointer-3}{}}
\newlabel{simplicial-lemma-cosk-fibre-product@cref}{{[originallemma][0][]018D}{[1][0][]source}{}{}{}}
\makeatother
\title{Base changes of hypercoverings along the gamma-covering condition}
\author{The Stacks Project Authors}
\date{stacks-project:hypercovering.tex:local-snapshot-20261012; GNU FDL 1.2 or later}
\begin{document}
\maketitle
\noindent\textit{Editorial note.} Statement span is lines 1510-1538 of sources/stacks-hypercovering/pinned-original-source.tex; proof span is lines 1540-1602, the complete author proof (61 body lines).\par\medskip
\noindent\textit{Editorial note.} Navigation stubs: 3 cited-but-not-included labels are supplied as non-load-bearing pointers with their pinned tags where the snapshot tag index has them.\par\medskip
\noindent\textit{Editorial note.} extra\_packages is empty on purpose: the retained preamble already loads xy with options and a second declaration clashes in the assembler wrapper.\par\medskip
\noindent\textit{Editorial note.} Text is retained verbatim, including the author's own typos and placeholders.\par\medskip
\noindent\textit{Editorial note.} Queue index 10 of queue\_hypercovering.json (pb=1823, stmt 1510-1538, proof 1540-1602).\par\medskip
\noindent\textit{Editorial note.} stacks-hypercovering wave-1 (hc01): fresh adjudication accept (2026-10-14); pool per standing fallback order (hypercovering -> pic -> spaces-more-cohomology) after md01 closed the moduli pool; hypercovering.tex 53 labels/34 with proof judged 26 accept / 8 reject / 14 pb=0 env-dup rows excluded; spans cite-free and reference/slogan/history-free so no patchenv needed.\par\medskip
\noindent\textit{Allowed standard prerequisites (editorial).}\begin{itemize}
\item Explicit prerequisite boundary. The cited inputs are prerequisites; the argument in the retained span is not replaced by any cited result.
\end{itemize}\medskip
\ifdefined\hypertarget\hypertarget{sourcepointer-1}{}\fi
\noindent\textit{Source reference.} Original lemma 01G5: sources/stacks-hypercovering/pinned-original-source.tex, source lines 563--575. Cited by the retained statement or proof but not part of the collected spans; supplied only as a navigation pointer. The argument does not depend on its content.\par\medskip
\ifdefined\hypertarget\hypertarget{sourcepointer-2}{}\fi
\noindent\textit{Source reference.} Original lemma 01G4: sources/stacks-hypercovering/pinned-original-source.tex, source lines 524--535. Cited by the retained statement or proof but not part of the collected spans; supplied only as a navigation pointer. The argument does not depend on its content.\par\medskip
\ifdefined\hypertarget\hypertarget{sourcepointer-3}{}\fi
\noindent\textit{Source reference.} Original lemma 018D: sources/stacks-hypercovering/pinned-original-source.tex, source lines 1602--1602. Cited by the retained statement or proof but not part of the collected spans; supplied only as a navigation pointer. The argument does not depend on its content.\par\medskip
\begin{lemma}
\label{lemma-funny-gamma}
Let $\mathcal{C}$ be a site with fibre products.
Let $X$ be an object of $\mathcal{C}$.
Let $K, L, M$ be simplicial objects of $\text{SR}(\mathcal{C}, X)$.
Let $a : K \to L$, $b : M \to L$ be morphisms.
Assume
\begin{enumerate}
\item $K$ is a hypercovering of $X$,
\item the morphism $M_0 \to L_0$ is a covering, and
\item for all $n \geq 0$ in the diagram
$$
\xymatrix{
M_{n + 1} \ar[dd] \ar[rr] \ar[rd]^\gamma &
&
(\text{cosk}_n \text{sk}_n M)_{n + 1} \ar[dd] \\
&
L_{n + 1}
\times_{(\text{cosk}_n \text{sk}_n L)_{n + 1}}
(\text{cosk}_n \text{sk}_n M)_{n + 1}
\ar[ld] \ar[ru]
& \\
L_{n + 1} \ar[rr] & & (\text{cosk}_n \text{sk}_n L)_{n + 1}
}
$$
the arrow $\gamma$ is a covering.
\end{enumerate}
Then the fibre product $K \times_L M$ is a hypercovering of $X$.
\end{lemma}

\begin{proof}
The morphism $(K \times_L M)_0 = K_0 \times_{L_0} M_0 \to K_0$
is a base change of a covering by (2), hence a covering, see
Lemma \ref{lemma-covering-permanence}. And $K_0 \to \{X \to X\}$
is a covering by (1). Thus $(K \times_L M)_0 \to \{X \to X\}$
is a covering by Lemma \ref{lemma-covering-permanence}. Hence
$K \times_L M$ satisfies the first condition of Definition
\ref{definition-hypercovering}.

\medskip\noindent
We still have to check that
$$
K_{n + 1} \times_{L_{n + 1}} M_{n + 1} = (K \times_L M)_{n + 1}
\longrightarrow
(\text{cosk}_n \text{sk}_n (K \times_L M))_{n + 1}
$$
is a covering for all $n \geq 0$. We abbreviate as follows:
$A = (\text{cosk}_n \text{sk}_n K)_{n + 1}$,
$B = (\text{cosk}_n \text{sk}_n L)_{n + 1}$, and
$C = (\text{cosk}_n \text{sk}_n M)_{n + 1}$.
The functor $\text{cosk}_n \text{sk}_n$ commutes with fibre products,
see Simplicial, Lemma \ref{simplicial-lemma-cosk-fibre-product}.
Thus the right hand side above is equal to $A \times_B C$.
Consider the following commutative diagram
$$
\xymatrix{
K_{n + 1} \times_{L_{n + 1}} M_{n + 1} \ar[r] \ar[d] &
M_{n + 1} \ar[d] \ar[rd]_\gamma \ar[rrd] &
& \\
K_{n + 1} \ar[r] \ar[rd] &
L_{n + 1} \ar[rrd] &
L_{n + 1} \times_B C \ar[l] \ar[r] &
C \ar[d] \\
&
A \ar[rr] &
&
B
}
$$
This diagram shows that
$$
K_{n + 1} \times_{L_{n + 1}} M_{n + 1}
=
(K_{n + 1} \times_B C)
\times_{(L_{n + 1} \times_B C), \gamma}
M_{n + 1}
$$
Now, $K_{n + 1} \times_B C \to A \times_B C$
is a base change of the covering $K_{n + 1} \to A$
via the morphism $A \times_B C \to A$, hence is a
covering. By assumption (3) the morphism $\gamma$ is a covering.
Hence the morphism
$$
(K_{n + 1} \times_B C)
\times_{(L_{n + 1} \times_B C), \gamma}
M_{n + 1}
\longrightarrow
K_{n + 1} \times_B C
$$
is a covering as a base change of a covering.
The lemma follows as a composition of coverings
is a covering.
\end{proof}

\end{document}
